% --- 1. Ajustes de Idioma y Texto ---
\usepackage[spanish, es-tabla, es-noshorthands]{babel}
\usepackage[utf8]{inputenc}
\usepackage[T1]{fontenc}
\usepackage{lmodern} % Fuente Latin Modern Vectorizada
\usepackage[nopatch=footnote]{microtype} % Optimiza el espaciado para evitar Underfull \hbox
\usepackage{ragged2e}
\usepackage{lipsum}
\usepackage{fancyhdr}
\usepackage{titlesec} % Para personalizar títulos de secciones
\usepackage{comment}
\usepackage{todonotes}

% --- 2. Geometría y Estilo de Página ---
\usepackage{geometry}
\geometry{
    letterpaper,
    top=2.5cm,
    bottom=2.5cm,
    inner=3cm,
    outer=3cm,
    headheight=20pt
}
\usepackage{float}
\usepackage{fancyhdr}
%\usepackage[hidelinks]{hyperref}
\usepackage{hyperref}

  \hypersetup{
    colorlinks=true,   % Activa el color en los elementos clickeables
    citecolor=blue,    % Pone los números de \cite{} en AZUL
    linkcolor=black,    % Pone las referencias internas (figuras, ecuaciones, secciones) en negro
    urlcolor=blue      % Pone los enlaces web en azul
}
\usepackage{parskip}
\usepackage{setspace}
\onehalfspacing % Activa el interlineado a 1.5 para todo el documento

% --- 3. Matemáticas y Física (Orden Crítico) ---
\usepackage{amsmath, amsfonts, amsthm, amssymb}
\usepackage{physics}
\usepackage{siunitx}
\AtBeginDocument{\RenewCommandCopy\qty\SI}
\usepackage{derivative}
\usepackage{gensymb} % Para \degree
\usepackage{textcomp} % Evita conflicto con gensymb (\perthousand)
\usepackage{mathrsfs} % Para fuentes matemáticas caligráficas
\numberwithin{equation}{chapter}

% --- 5. Tablas y Figuras ---
\usepackage{graphicx, wrapfig}
\usepackage{tabularx, multirow, multicol, array, booktabs, bigstrut}
\usepackage{caption, subcaption}
\usepackage{xcolor, colortbl}

% --- 6. Misceláneos ---
\usepackage{comment, tcolorbox, schemata}

% --- 7. Ajustes Globales de Formato ---
\usepackage[bottom]{footmisc} % Obliga a las notas al pie a fijarse al fondo de la página

\sloppy % Relaja la justificación para eliminar los molestos Underfull \hbox
\overfullrule=0pt % Cambia a 5pt si quieres ver dónde se sale el texto del margen

\raggedbottom                 % Absorbe el espacio en blanco sobrante sin estirar el texto
\setcounter{secnumdepth}{3} % Niveles en el texto
\setcounter{tocdepth}{3} % Niveles en el índice

\usepackage{csquotes}
\usepackage[
    backend=biber,
    style=numeric-comp,  % Citas numeradas [1] y agrupadas [1-3]
    sorting=none,        % Orden de aparición en el texto [1], [2], [3]...
    giveninits=true,     % Usa iniciales en los nombres (ej. M. D. Rotondaro)
    terseinits=false,    % Mantiene puntos en las iniciales (M. D.)
    %articletitle=true,   % Muestra siempre el título completo de los artículos
    doi=true,            % Muestra enlaces DOI activos
    isbn=false,          % Oculta el ISBN para limpiar la vista
    url=false            % Oculta URLs redundantes si ya existe un DOI
]{biblatex}

% Añade espacio vertical entre cada referencia (para que respire y no se vea apretada)
\setlength{\bibitemsep}{0.8em}

\titleformat{\chapter}[display]
  {\normalfont\huge\bfseries}{\chaptertitlename\ \thechapter}{10pt}{\Huge}
\titleformat{\paragraph}
  {\normalfont\normalsize\bfseries} % Texto normal, negrita
  {\theparagraph}                   % El número (ej. 1.1.1.1)
  {1em}                             % Espacio tras el número
  {}                                % Código extra

\titlespacing*{\chapter}{0pt}{80pt}{10pt}% Separación del título de capítulo (izquierda, superior, inferior)
\titlespacing*{\section}{0pt}{10pt}{0pt}% Separación del título de subsección (izquierda, superior, inferior)
\titlespacing*{\subsection}{0pt}{10pt}{0pt}% Separación del título de subsección (izquierda, superior, inferior)
\titlespacing*{\paragraph}{0pt}{0pt}{0pt}

% --- Configuración de Encabezados y Pies de Página ---
\pagestyle{fancy}
\fancyhf{} % Limpia todos los campos de encabezado y pie de página

% --- CONTROL DE MAYÚSCULAS/MINÚSCULAS EN ENCABEZADOS ---
% 1. Capítulo: fuerza MAYÚSCULAS (\MakeUppercase)
\renewcommand{\chaptermark}[1]{%
  \markboth{\MakeUppercase{\chaptername\ \thechapter.\ #1}}{}%
}

% 2. Sección: respeta el texto tal cual lo escribiste (Formato oración)
\renewcommand{\sectionmark}[1]{%
  \markright{\thesection.\ #1}%
}

% Páginas impares (Odd - O)
\fancyhead[LO]{\rightmark} % Izquierda (L): Subsección
\fancyhead[RO]{\thepage}   % Derecha (R): Número de página

% Páginas pares (Even - E)
\fancyhead[LE]{\thepage}   % Izquierda (L): Número de página
\fancyhead[RE]{\leftmark}  % Derecha (R): Sección

% Línea divisoria del encabezado (opcional)
\renewcommand{\headrulewidth}{0.4pt}

% ENCABEZADO LIMPIO PARA INICIO DE CAPÍTULO
% En 'report', los capítulos usan automáticamente el estilo 'plain'.
\assignpagestyle{\chapter}{plain}
\fancypagestyle{plain}{%
  \fancyhf{}                          % Limpia encabezado y pie de página
  \fancyhead[RO,LE]{\thepage}         % Muestra SOLO el número arriba (derecha impares, izquierda pares)
  \renewcommand{\headrulewidth}{0pt}  % Elimina la línea del encabezado
  \renewcommand{\footrulewidth}{0pt}  % Elimina la línea del pie
}

